% Section 27: The residue theorem as Stokes' theorem} | Part: 6 | Depends on: see docs/REWRITE_PLAN.md | Status: todo | Original lines: 3716-3723
%!TEX root = ../../../main.tex
% =====================================================================
\section{The residue theorem as Stokes' theorem}\label{sec:complex}
% =====================================================================

Is the residue theorem related to all of this? Yes: it is not a separate theorem but Stokes' theorem applied to a particular kind of complex-valued $1$-form on a region with holes.
Cauchy's theorem is Stokes for a closed form; the residue theorem is Stokes for a closed form that fails to be exact exactly at the punctures, which is the mechanism of Example~\ref{ex:angle}
(the angle form $\dd\theta$ on $\R^2\setminus\{0\}$). Its real and imaginary parts are, respectively, the circulation and the flux of Sections~\ref{sec:stokes}--\ref{sec:phys}.
% Section 27: Complex-valued $1$-forms and the contraction} | Part: 6 | Depends on: see docs/REWRITE_PLAN.md | Status: todo | Original lines: 3724-3743
%!TEX root = ../../../main.tex
\subsection{Complex-valued $1$-forms and the contraction}\label{sec:complexforms}

Identify $\C=\R^2$ with $z=x+iy$. The form $\dd z=\dd x+i\,\dd y$ is a complex-valued $1$-form (a $\C$-valued covector at each point), and its contraction with a tangent vector $v=(v^x,v^y)\simeq v^x+iv^y$ is simply $\dd z(v)=v$. For a function $f:U\subset\C\to\C$ the form $f\,\dd z$ contracts with a velocity as $\langle f\,\dd z,\dot\gamma(t)\rangle=f(\gamma(t))\dot\gamma(t)$,
\begin{equation}\label{eq:cint}
\int_\gamma f\,\dd z=\int_a^bf(\gamma(t))\,\dot\gamma(t)\,\dd t,
\end{equation}
which is Definition~\eqref{eq:lineint} extended $\C$-linearly. For every smooth $f$ one has $\dd f=\partial_zf\,\dd z+\partial_{\bar z}f\,\dd\bar z$, and
\begin{equation}\label{eq:holo}
f\text{ holomorphic}\iff\partial_{\bar z}f=0\iff\dd f=f'(z)\,\dd z\iff\dd f_p:\R^2\to\C\text{ is }\C\text{-linear}.
\end{equation}
The Cauchy--Riemann equations are thus the statement that the covector $\dd f$ is complex-linear: $\dd f(iv)=i\,\dd f(v)$.

\emph{Steps behind these statements.} With $\partial_z=\frac12(\partial_x-i\partial_y)$ and $\partial_{\bar z}=\frac12(\partial_x+i\partial_y)$:
\begin{align*}
\partial_zf\,\dd z+\partial_{\bar z}f\,\dd\bar z&=\tfrac12(f_x-if_y)(\dd x+i\,\dd y)+\tfrac12(f_x+if_y)(\dd x-i\,\dd y)\\
&=\tfrac12\bigl(f_x\dd x+if_x\dd y-if_y\dd x+f_y\dd y\bigr)+\tfrac12\bigl(f_x\dd x-if_x\dd y+if_y\dd x+f_y\dd y\bigr)\\
&=f_x\dd x+f_y\dd y=\dd f.
\end{align*}
Writing $f=u+iv$, the condition $\partial_{\bar z}f=0$ means $f_x+if_y=(u_x-v_y)+i(v_x+u_y)=0$, \ie the usual Cauchy--Riemann equations $u_x=v_y$, $u_y=-v_x$. Finally $\dd f=f'\dd z$ means $\partial_{\bar z}f=0$ and $\partial_zf=f'$.

% Section 27.2: Cauchy's theorem and the Cauchy--Riemann equations as ``$\operatorname{div}=0$, $\operatorname{curl}=0$'' | Part: 6 | Depends on: sec:stokes, sec:ga-calculus | Status: todo | Original lines: 3744-3778
%!TEX root = ../../../main.tex
\subsection{Cauchy's theorem and the Cauchy--Riemann equations as ``$\operatorname{div}=0$, $\operatorname{curl}=0$''}\label{sec:cauchy}

Let us compute $\dd(f\,\dd z)$ in three steps.
\begin{enumerate}
\item Leibniz rule and $\dd(\dd z)=0$: $\dd(f\,\dd z)=\dd f\wedge\dd z$.
\item Insert $\dd f=\partial_zf\,\dd z+\partial_{\bar z}f\,\dd\bar z$ and use $\dd z\wedge\dd z=0$: $\dd f\wedge\dd z=\partial_{\bar z}f\,\dd\bar z\wedge\dd z$.
\item Expand $\dd\bar z\wedge\dd z=(\dd x-i\,\dd y)\wedge(\dd x+i\,\dd y)=i\,\dd x\wedge\dd y-i\,\dd y\wedge\dd x=2i\,\dd x\wedge\dd y$.
\end{enumerate}
\begin{equation}\label{eq:dfdz}
\dd(f\,\dd z)=\partial_{\bar z}f\,\dd\bar z\wedge\dd z=2i\,\partial_{\bar z}f\,\dd x\wedge\dd y .
\end{equation}
Hence $f$ is holomorphic if and only if the $1$-form $f\,\dd z$ is closed, and Stokes' theorem gives at once

\begin{theorem}[Cauchy]\label{thm:cauchy}
If $f$ is holomorphic on a neighbourhood of a compact region $\Omega\subset\C$ with piecewise smooth boundary, then
$\displaystyle\oint_{\partial\Omega}f\,\dd z=\int_\Omega\dd(f\,\dd z)=0.$
\end{theorem}

\subhead{Real and imaginary parts: circulation and flux}
Write $w=u-iv$ (the ``complex velocity'' of a planar vector field $\mathbf u=(u,v)$). Then
\begin{equation}\label{eq:wdz}
w\,\dd z=\underbrace{(u\,\dd x+v\,\dd y)}_{\mathbf u^\flat}+i\underbrace{(u\,\dd y-v\,\dd x)}_{\iota_{\mathbf u}(\dd x\wedge\dd y)=\hodge\mathbf u^\flat}.
\end{equation}
The real part is the covector $\mathbf u^\flat$ (its integral is the circulation $\Gamma$), and the imaginary part is the contraction $\iota_{\mathbf u}\vol$ (its integral is the flux $Q$ through the curve, cf.\ \eqref{eq:flux}). The split follows from $w\,\dd z=(u-iv)(\dd x+i\,\dd y)=u\,\dd x+iu\,\dd y-iv\,\dd x+v\,\dd y$ and from $\iota_{\mathbf u}(\dd x\wedge\dd y)=u\,\dd y-v\,\dd x$ (Example~\ref{ex:interiorplane}). Moreover, from \eqref{eq:dfdz}, in steps:
\[
\partial_{\bar z}w=\tfrac12(\partial_x+i\partial_y)(u-iv)=\tfrac12\bigl[(u_x+v_y)+i(u_y-v_x)\bigr]=\tfrac12\bigl(\operatorname{div}\mathbf u-i\operatorname{curl}\mathbf u\bigr)
\]
(with $\operatorname{curl}\mathbf u=\partial_xv-\partial_yu$), 
so \emph{$w$ is holomorphic $\iff$ the flow is incompressible ($\operatorname{div}\mathbf u=0$) and irrotational ($\operatorname{curl}\mathbf u=0$).}
Cauchy's theorem is then Green's theorem (Kelvin--Stokes in the plane) for the real part and the divergence theorem (Gauss) for the imaginary part, both applied at the same time.

\begin{specialcase}[the Hodge star in dimension two]\label{sc:hodge2}
In the language of Section~\ref{sec:classical}, in dimension two the Hodge star on $1$-forms satisfies $\hodge\dd z=-i\,\dd z$ and $\hodge\dd\bar z=i\,\dd\bar z$: $\dd z$ is an eigenform of $\hodge$, and since $\hodge$ on $1$-forms depends only on the conformal class of the metric, holomorphy is a conformally invariant notion (conformal maps preserve holomorphic functions). The eigenform property is checked in one line: $\hodge\dd x=\dd y$, $\hodge\dd y=-\dd x$, so $\hodge\dd z=\dd y-i\,\dd x=-i(\dd x+i\,\dd y)=-i\,\dd z$.
\end{specialcase}

% Section 27.3: The residue theorem | Part: 6 | Depends on: sec:stokes, sec:ga-calculus | Status: todo | Original lines: 3779-3844
%!TEX root = ../../../main.tex
\subsection{The residue theorem}\label{sec:residue}

Let $f$ be holomorphic on an open set containing $\overline\Omega$ except at finitely many points $a_1,\dots,a_k\in\Omega$ where it has isolated singularities (poles or essential singularities). The \emph{residue} of $f$ at $a$ is the coefficient $c_{-1}$ of its Laurent expansion $f(z)=\sum_nc_n(z-a)^n$, equivalently
\begin{equation}\label{eq:residue}
\Res_af=\frac1{2\pi i}\oint_{|z-a|=\varepsilon}f\,\dd z\qquad(\text{small }\varepsilon>0,\ \text{counter-clockwise}).
\end{equation}
\emph{Why the two definitions agree (steps).} Parametrise the circle: $z=a+\varepsilon e^{it}$, $t\in[0,2\pi]$, so $\dd z=i\varepsilon e^{it}\,\dd t$. For an integer $n$:
\[
\oint(z-a)^n\,\dd z=\int_0^{2\pi}\varepsilon^ne^{int}\,i\varepsilon e^{it}\,\dd t=i\varepsilon^{n+1}\int_0^{2\pi}e^{i(n+1)t}\,\dd t=\begin{cases}2\pi i,&n=-1,\\0,&n\neq-1,\end{cases}
\]
because for $m=n+1\neq0$, $\int_0^{2\pi}e^{imt}\,\dd t=\bigl[e^{imt}/(im)\bigr]_0^{2\pi}=0$, while for $m=0$ the integrand is $1$. Integrating the Laurent series term by term (it converges uniformly on the circle) leaves only $c_{-1}\cdot2\pi i$.

% Figure | used in parts/06-complex-analysis/01-residue-theorem/03-residue-theorem.tex
%!TEX root = ../../main.tex
\begin{figure}[ht]
\centering
\begin{tikzpicture}[>=Stealth,scale=1.1]
  \fill[blue!12] (0,0) circle (2.0);
  \draw[navy] (0,0) circle (2.0);
  \foreach \c/\n in {(-0.9,0.55)/1,(0.95,0.75)/2,(0.35,-0.9)/3}{
     \fill[white] \c circle (0.3);
     \draw[navy] \c circle (0.3);
     \fill \c circle (0.03);
  }
  \node[navy,left] at (-0.9,0.55) {\small$a_1$};
  \node[navy,right] at (0.95,0.75) {\small$a_2$};
  \node[navy,below] at (0.35,-0.9) {\small$a_3$};
  \node[navy] at (-0.3,-0.1) {$\Omega_\varepsilon$};
  \node[navy,above right] at (1.5,1.45) {$\partial\Omega$};
  \draw[->,navy,thick] (225:2.0) arc (225:240:2.0);
  \draw[->,navy,thick] ($(-0.9,0.55)+(30:0.3)$) arc (30:-20:0.3);
  \draw[->,navy,thick] ($(0.95,0.75)+(150:0.3)$) arc (150:200:0.3);
  \draw[->,navy,thick] ($(0.35,-0.9)+(60:0.3)$) arc (60:10:0.3);
\end{tikzpicture}
\caption{Proof of the residue theorem. Stokes on $\Omega_\varepsilon=\Omega$ minus small discs: the outer boundary runs counter-clockwise and the small circles clockwise (region on the left). Since $f\,\dd z$ is closed on $\Omega_\varepsilon$, the integral over the outer boundary equals the sum of the integrals over the small circles traversed counter-clockwise. This is the annulus of Figure~\ref{fig:annulus} with several holes.}
\label{fig:residue}
\end{figure}


\begin{theorem}[Residue theorem]\label{thm:residue}
Under these assumptions,
\[
\oint_{\partial\Omega}f\,\dd z=2\pi i\sum_{j=1}^k\Res_{a_j}f .
\]
\end{theorem}
\begin{proof}
Remove from $\Omega$ small open discs $D_j=D(a_j,\varepsilon)$ and let $\Omega_\varepsilon=\Omega\setminus\bigcup_jD_j$, a manifold with boundary $\partial\Omega_\varepsilon=\partial\Omega-\sum_j\partial D_j$ (the circles are traversed clockwise: induced orientation, Figure~\ref{fig:residue}). On $\Omega_\varepsilon$ the form $f\,\dd z$ is smooth and closed, so Stokes' theorem \eqref{eq:stokes} gives
\[
0=\int_{\Omega_\varepsilon}\dd(f\,\dd z)=\oint_{\partial\Omega}f\,\dd z-\sum_j\oint_{\partial D_j^{\mathrm{ccw}}}f\,\dd z,
\]
and each small-circle integral is $2\pi i\Res_{a_j}f$ by \eqref{eq:residue}.
\end{proof}

\begin{example}[Two applications of the residue theorem]\label{ex:residues}
(a) $\displaystyle\oint_{|z|=2}\frac{\dd z}{z(z-1)}$. Partial fractions: $\frac1{z(z-1)}=\frac1{z-1}-\frac1z$, so $\Res_1=+1$ and $\Res_0=-1$ (both poles lie inside $|z|<2$). The theorem gives $2\pi i(1-1)=0$.
\emph{Check:} $\oint\frac{\dd z}{z-1}=2\pi i$ and $\oint\frac{\dd z}z=2\pi i$ by the circle computation above, and their difference is $0$.

(b) $\displaystyle\oint_{|z|=1}\frac{e^z}{z^2}\,\dd z$. Expand $e^z=1+z+\frac{z^2}2+\cdots$, so $\frac{e^z}{z^2}=\frac1{z^2}+\frac1z+\frac12+\cdots$ and $c_{-1}=1$. The integral is $2\pi i\cdot1=2\pi i$. (Equivalently, $\Res_0=\frac{\dd}{\dd z}e^z\big|_0=1$, the formula for a double pole.)
\end{example}

\begin{keyidea}[why only $c_{-1}$ survives: exact versus closed]
On $\C\setminus\{a\}$ the forms $(z-a)^n\dd z$ with $n\neq-1$ are \emph{exact}: $(z-a)^n\dd z=\dd\bigl(\frac{(z-a)^{n+1}}{n+1}\bigr)$, so they integrate to zero on any closed loop (Section~\ref{sec:expl-closed}). The only closed, non-exact one is
\begin{equation}\label{eq:dzoverz}
\frac{\dd z}{z-a}=\dd\ln|z-a|+i\,\dd\theta,
\end{equation}
whose imaginary part is precisely the angle form $\dd\theta$ of Example~\ref{ex:angle}, with $\oint\dd\theta=2\pi$; hence $\oint\frac{\dd z}{z-a}=i\cdot2\pi=2\pi i$. Derivation of the identity: write $z=re^{i\theta}$; then $\dd z=e^{i\theta}\dd r+ire^{i\theta}\dd\theta$ and $\frac{\dd z}z=\frac{\dd r}r+i\,\dd\theta=\dd\ln r+i\,\dd\theta$.
In other words, the residue measures the component of $f\,\dd z$ along the one cohomology class which is not exact near $a$.
\end{keyidea}

% Section 27: The cohomological meaning of the residue} | Part: 6 | Depends on: see docs/REWRITE_PLAN.md | Status: todo | Original lines: 3845-3854
%!TEX root = ../../../main.tex
\subsection{The cohomological meaning of the residue}\label{sec:rescohom}

For $X=\C\setminus\{a_1,\dots,a_k\}$ one has $H^1_{\mathrm{dR}}(X;\C)\cong\C^k$, generated by the classes of $\dd z/(z-a_j)$, and $H_1(X;\Z)\cong\Z^k$, generated by small circles around the $a_j$. The residue is the de Rham pairing between these:
$\Res_{a_j}f=\frac1{2\pi i}\bigl\langle[\text{circle around }a_j],[f\,\dd z]\bigr\rangle$.
A closed loop $\gamma$ in $X$ is homologous to $\sum_jn_j[\text{circle}_j]$, where $n_j=\frac1{2\pi i}\oint_\gamma\frac{\dd z}{z-a_j}$ is the winding number of $\gamma$ around $a_j$ (Example~\ref{ex:angle}); \eg a loop that goes twice around $a_j$ has $\oint\frac{\dd z}{z-a_j}=2\cdot2\pi i$, so $n_j=2$. This yields the general version
\begin{equation}\label{eq:winding}
\oint_\gamma f\,\dd z=2\pi i\sum_jn(\gamma,a_j)\Res_{a_j}f .
\end{equation}
The integral depends only on the homology class of the contour: a contour can be deformed freely as long as it does not cross a singularity, which is the practical principle behind contour integration.

% Section 27: Beyond holomorphic functions: Cauchy--Pompeiu and the Coulomb law in 2D} | Part: 6 | Depends on: see docs/REWRITE_PLAN.md | Status: todo | Original lines: 3855-3873
%!TEX root = ../../../main.tex
\subsection{Beyond holomorphic functions: Cauchy--Pompeiu and the Coulomb law in 2D}\label{sec:pompeiu}

If $f$ is merely smooth, Stokes still applies. Let us derive the Cauchy--Pompeiu formula step by step, with $\omega=\dfrac{f\,\dd z}{z-z_0}$ on $\Omega_\varepsilon=\Omega\setminus D(z_0,\varepsilon)$.
\begin{enumerate}
\item \emph{Differentiate.} $1/(z-z_0)$ is holomorphic away from $z_0$, so $\partial_{\bar z}\bigl(f/(z-z_0)\bigr)=\partial_{\bar z}f/(z-z_0)$ and, by \eqref{eq:dfdz}, $\dd\omega=\dfrac{\partial_{\bar z}f}{z-z_0}\,2i\,\dd x\wedge\dd y$.
\item \emph{Stokes on $\Omega_\varepsilon$} (outer boundary ccw, small circle cw): $\displaystyle\oint_{\partial\Omega}\omega-\oint_{|z-z_0|=\varepsilon}\omega=\int_{\Omega_\varepsilon}\dd\omega$.
\item \emph{The small circle.} With $z=z_0+\varepsilon e^{it}$: $\displaystyle\oint\omega=\int_0^{2\pi}\frac{f(z_0+\varepsilon e^{it})}{\varepsilon e^{it}}\,i\varepsilon e^{it}\,\dd t=i\int_0^{2\pi}f(z_0+\varepsilon e^{it})\,\dd t\xrightarrow[\varepsilon\to0]{}2\pi if(z_0)$ by continuity of $f$.
\item \emph{Let $\varepsilon\to0$.} The area integral converges because $1/|z-z_0|$ is integrable in the plane (area element $r\,\dd r\,\dd\theta$). We get $\displaystyle\oint_{\partial\Omega}\frac{f\,\dd z}{z-z_0}-2\pi if(z_0)=2i\iint_\Omega\frac{\partial_{\bar z}f}{z-z_0}\,\dd A$.
\item \emph{Divide by $2\pi i$ and solve for $f(z_0)$:}
\begin{equation}\label{eq:pompeiu}
f(z_0)=\frac1{2\pi i}\oint_{\partial\Omega}\frac{f(z)}{z-z_0}\,\dd z-\frac1\pi\iint_\Omega\frac{\partial_{\bar z}f}{z-z_0}\,\dd A .
\end{equation}
\end{enumerate}
For holomorphic $f$ the area term vanishes and we recover Cauchy's integral formula. In the language of distributions the same computation says $\partial_{\bar z}\bigl(\frac1{\pi z}\bigr)=\delta_0$, \ie $1/(\pi z)$ is the fundamental solution of $\partial_{\bar z}$, in perfect analogy with the Gauss law in the plane (Section~\ref{sec:phys-em}): a point charge $Q$ (or a point vortex with circulation $\Gamma$ and source $Q$) is a simple pole of the complex field $w$ whose residue encodes the charge: $\Res_0w=\frac{\Gamma+iQ}{2\pi i}$, \ie $2\pi i\Res_0w=\Gamma+iQ$,
\[
w(z)=\frac{\Gamma+iQ}{2\pi i}\,\frac1z\quad\Longrightarrow\quad\oint w\,\dd z=\Gamma+iQ .
\]
This is why the residue theorem reads, in electrostatics or fluid mechanics, as ``the total flux (and circulation) through a curve equals the sum of the charges (vortices) inside''.

% Section 27: The argument principle: counting zeros with a $1$-form} | Part: 6 | Depends on: see docs/REWRITE_PLAN.md | Status: todo | Original lines: 3874-3897
%!TEX root = ../../../main.tex
\subsection{The argument principle: counting zeros with a $1$-form}\label{sec:argument}

Let $f$ be meromorphic on a neighbourhood of the closure of $\Omega$, with no zeros or poles on $\partial\Omega$. Then $\frac{f'}f\dd z=\dd\log f=\dd\ln|f|+i\,\dd\arg f=f\pb\bigl(\frac{\dd z}z\bigr)$. The first term is exact, so the residue theorem becomes
\begin{equation}\label{eq:argument}
\frac1{2\pi i}\oint_{\partial\Omega}\frac{f'}f\,\dd z=\frac1{2\pi}\oint_{\partial\Omega}f\pb\dd\theta=N-P,
\end{equation}
where $N$ and $P$ are the numbers of zeros and poles in $\Omega$ counted with multiplicity: the winding number of the curve $f(\partial\Omega)$ around $0$.
\emph{Steps.} (1) Locally $\log f=\ln|f|+i\arg f$, and $\dd\log f=\frac{\dd f}f=\frac{f'}f\dd z$; this is also the pull-back $f\pb(\dd z/z)$ because $f\pb\dd z=\dd f=f'\dd z$. (2) $\dd\ln|f|$ is exact, so its integral over a closed curve is $0$. (3) What remains is $i\oint\dd\arg f=2\pi i\cdot(\text{winding number of }f\circ\gamma\text{ around }0)$.
(4) By the residue theorem, $\frac{f'}f$ has a simple pole with residue $+m$ at a zero of order $m$ and $-m$ at a pole of order $m$ (since $f=(z-a)^{\pm m}g$ gives $f'/f=\pm m/(z-a)+g'/g$).
\numcheck $f(z)=z^2$ on $|z|=1$: $\frac{f'}f=\frac2z$, $\frac1{2\pi i}\oint\frac2z\,\dd z=2=N$ (a double zero at the origin), and $f(e^{it})=e^{2it}$ winds twice around $0$.
This is the degree of the map $f|_{\partial\Omega}:\partial\Omega\to\C\setminus\{0\}$, computed by pulling back the angle form (Section~\ref{sec:stokes}), and it is the basis of Rouch\'e's theorem and of the fundamental theorem of algebra.

\begin{example}[A real integral by residues, step by step]\label{ex:realint}
To compute $I=\displaystyle\int_{-\infty}^\infty\frac{\dd x}{1+x^2}$, apply Theorem~\ref{thm:residue} to $f\,\dd z$ with $f=\dfrac1{1+z^2}$ on the upper half-disc $\Omega_R$ of radius $R>1$, whose boundary is the segment $[-R,R]$ followed by the semicircle $C_R$.
\begin{enumerate}
\item \emph{Poles.} $1+z^2=(z-i)(z+i)$; only $z=i$ lies inside $\Omega_R$.
\item \emph{Residue.} $\Res_if=\lim_{z\to i}\frac{z-i}{(z-i)(z+i)}=\frac1{2i}$.
\item \emph{Theorem.} $\int_{-R}^Rf(x)\,\dd x+\int_{C_R}f\,\dd z=2\pi i\cdot\frac1{2i}=\pi$. (On the real axis the pull-back of $f\,\dd z$ is $f(x)\,\dd x$, the contraction with the velocity $1$.)
\item \emph{The arc vanishes.} On $|z|=R$, $|1+z^2|\ge R^2-1$, so $\bigl|\int_{C_R}f\,\dd z\bigr|\le\dfrac{\pi R}{R^2-1}\to0$ as $R\to\infty$ (length $\pi R$ times the bound on $|f|$).
\item \emph{Limit.} Letting $R\to\infty$: $I=\pi$.
\end{enumerate}
\emph{Check:} $\int\frac{\dd x}{1+x^2}=\arctan x\big|_{-\infty}^\infty=\frac\pi2-(-\frac\pi2)=\pi$ \checkmark.
\end{example}

% Section 27: The current picture} | Part: 6 | Depends on: see docs/REWRITE_PLAN.md | Status: todo | Original lines: 3898-3923
%!TEX root = ../../../main.tex
\subsection{The current picture}\label{sec:rescurrent}

\begin{itemize}
\item \emph{Riemann surfaces.} If $\eta$ is a meromorphic $1$-form on a compact Riemann surface $S$ (no boundary), Stokes with small discs removed gives $\sum_{p\in S}\Res_p\eta=0$: the sum of all residues vanishes. Combined with Poincar\'e/Serre duality this is a key ingredient in the proofs of the Riemann--Roch theorem and in the theory of abelian differentials. Note that this statement does not even need a metric: $\eta$ is a $1$-form and $\Res_p\eta$ is intrinsic, which is why residues are defined for forms $f\,\dd z$ and not for functions $f$ (the residue of the function $f$ changes under a change of coordinate $z\mapsto h(z)$, that of the form $f\,\dd z$ does not).
\item \emph{Several complex variables.} The kernel $1/(z-z_0)$ is replaced by the Bochner--Martinelli kernel, and residues by Poincar\'e--Leray residues of closed forms with poles along hypersurfaces; both are proved with Stokes' theorem and are expressed through de Rham cohomology of the complement of the singular set.
\item \emph{Currents.} Formulae such as $\partial_{\bar z}(1/\pi z)=\delta_0$ and the Poincar\'e--Lelong formula $\frac i\pi\partial\bar\partial\log|f|=[\{f=0\}]$ treat the residue as the current of integration over the singular set (Section~\ref{sec:current}, item~3), exactly as Stokes' theorem is treated in geometric measure theory.
\end{itemize}

\begin{center}
\small
\renewcommand{\arraystretch}{1.3}
\begin{tabular}{@{}p{4.6cm}p{9.6cm}@{}}
\toprule
\textbf{Complex analysis} & \textbf{Differential forms}\\
\midrule
$f$ holomorphic & $f\,\dd z$ closed; $\dd f$ is $\C$-linear\\
Cauchy's theorem & Stokes for a closed form on $\Omega$\\
Cauchy--Riemann equations & $\operatorname{div}=0$ and $\operatorname{curl}=0$ (real/imaginary parts of $w\,\dd z$)\\
Residue $\Res_af$ & period of $f\,\dd z$ over a small circle $=$ de Rham pairing\\
Residue theorem & Stokes on $\Omega$ minus discs; boundary $=$ outer curve $-$ small circles\\
Winding number, argument principle & $\frac1{2\pi}\oint f\pb\dd\theta=$ degree of a map to $S^1$\\
Cauchy--Pompeiu, $\partial_{\bar z}(1/\pi z)=\delta$ & Stokes with $\dd\omega\neq0$; Gauss law in the plane\\
\bottomrule
\end{tabular}
\end{center}

% Section 27.8: Synthesis IV: complex analysis in the geometric algebra of the plane | Part: 6 | Depends on: sec:stokes, sec:ga-calculus | Status: todo | Original lines: 3924-3950
%!TEX root = ../../../main.tex
\subsection{Synthesis IV: complex analysis in the geometric algebra of the plane}\label{sec:synth4}

\begin{synthesis}[where does the imaginary unit come from?]
The symbol $i$ has appeared in three guises. (1) As a scalar: $\C=\Cl(0,1)$, the geometric algebra of a line with $e^2=-1$ (Section~\ref{sec:clifford}). (2) As a bivector: in the plane algebra $\Cl(2,0)$ the unit area element $\pss=e_1e_2$ satisfies $\pss^2=-1$, and the even subalgebra $\{x+\pss y\}$ is $\C$; multiplying a vector on the left by $e_1$ identifies $\R^2$ with $\C$, $e_1(x+\pss y)=xe_1+ye_2$. (3) As the factor in complex-valued forms such as $\dd z=\dd x+i\,\dd y$, where $i$ is a formal scalar commuting with everything.
In option (2) $i$ is the \emph{oriented plane itself}, and ``multiplication by $i$'' is a quarter-turn \emph{because} $\pss$ rotates vectors by $90^\circ$ under the geometric product: the dual of a vector in $2$D is the vector turned by a quarter turn (Section~\ref{sec:duality}).
\end{synthesis}

\begin{proposition}[Cauchy--Riemann as a single multivector equation]\label{prop:CRga}
Let $\mathbf u=u\,e_1+v\,e_2$ be a vector field in the plane, $\vd=e_1\partial_x+e_2\partial_y$, $\pss=e_1e_2$, and $w=u-iv$ the complex velocity of \eqref{eq:wdz}. Then
\[
\vd\,\mathbf u=\operatorname{div}\mathbf u+\pss\,\operatorname{curl}\mathbf u=2\,\partial_z\bar w,
\]
so $w$ is holomorphic if and only if $\vd\mathbf u=0$ (a \emph{monogenic} function).
\end{proposition}
\begin{proof}
$\vd\mathbf u=(e_1\partial_x+e_2\partial_y)(ue_1+ve_2)=(u_x+v_y)+e_1e_2(v_x-u_y)$, using $e_1e_2=-e_2e_1$ and $e_i^2=1$; this is $\operatorname{div}\mathbf u+\pss\operatorname{curl}\mathbf u$. On the other hand $2\partial_z\bar w=(\partial_x-i\partial_y)(u+iv)=(u_x+v_y)+i(v_x-u_y)$ with $i\leftrightarrow\pss$. Holomorphy of $w$ means $\partial_{\bar z}w=0$, whose conjugate is $\partial_z\bar w=0$.
\end{proof}

\begin{example}[a monogenic field]\label{ex:monogenic}
$w=z^2=(x^2-y^2)+2ixy$. Since $w=u-iv$, we read off $u=x^2-y^2$ and $v=-2xy$. Then $\operatorname{div}\mathbf u=2x-2x=0$ and $\operatorname{curl}\mathbf u=v_x-u_y=-2y+2y=0$, so $\vd\mathbf u=0$ \checkmark.
\end{example}

\begin{synthesis}[Cauchy's theorem, three ways]
\emph{Forms:} $\dd(f\,\dd z)=2i\,\partial_{\bar z}f\,\dd x\wedge\dd y$ (equation \eqref{eq:dfdz}) and Stokes. \emph{Vector calculus:} Green's theorem for the real part and the divergence theorem for the imaginary part of \eqref{eq:wdz}. \emph{Geometric algebra:} the single equation $\vd\mathbf u=0$ of Proposition~\ref{prop:CRga}; in higher dimension the same equation $\vd F=0$ defines monogenic functions (Clifford analysis), for which an analogue of Cauchy's integral formula holds (Hestenes--Sobczyk, Section~\ref{sec:refs}).
\end{synthesis}

